\section{Related Work}

\textbf{Trustworthy auditing and data governance.}
Trustworthy AI evaluation increasingly requires evidence about deployed behavior rather than only benchmark accuracy, and foundation-model deletion is a hard audit target because residual dependence can appear through retrieval, adaptation, memory, or generated derivatives. Recent challenge work on federated foundation models lists private-data utilization, continual learning, unlearning, and watermarking among the major open problems~\cite{fan2025fedfmchallenges}. \method{} contributes to this agenda by making deletion claims testable through declared probes, confidence bounds, and artifact-level provenance; unlike explainability or attribution work, it asks whether a specific deleted source still contributes to post-deletion behavior, not merely which context was salient for one answer.

\textbf{Machine unlearning.}
Machine unlearning removes the effect of selected training data after deployment~\cite{cao2015machineunlearning,bourtoule2021machineunlearning,ginart2019making}, and foundation-model unlearning extends this to large pretrained and adapted models where exact retraining is infeasible~\cite{yao2024largeunlearning,maini2024tofu,liu2024rethinking,wang2024llmunlearning}. Benchmarks such as TOFU, MUSE, RWKU, and WMDP evaluate fictitious-biography unlearning, six-way unlearning properties, real-world knowledge removal, and hazardous-knowledge reduction~\cite{maini2024tofu,shi2024muse,jin2024rwku,li2024wmdp}, and surveys emphasize evaluating both forgetting and retained utility~\cite{jin2024surveyunlearning}. We differ by focusing on auditability across the data pipeline rather than proposing a new weight-update rule: \method{} can wrap a TOFU-, MUSE-, RWKU-, or WMDP-style run and ask whether downstream artifacts, caches, derivatives, or adapters remain behaviorally distinguishable from a clean reference pipeline.

\textbf{Retrieval-augmented generation and provenance.}
RAG conditions generation on retrieved documents~\cite{lewis2020rag,karpukhin2020dpr,izacard2021fid,gao2024ragsurvey}, which improves grounding but introduces deletion risk: removing a source from one index need not remove derived chunks, summaries, cached contexts, or fine-tuning examples, and multi-stage pipelines can have a clean first-stage index while a reranker, cache, or derivative reintroduces a deleted fact. Work on attribution, context usage, source tracing, and RAG evaluation~\cite{gao2023retrievalaugmented,asai2024selfrag,yu2024ragevalsurvey,akkiraju2024facts} treats retrieval quality as the object of study; \method{} instead treats these artifacts as first-class audit objects and adds deletion risk as an orthogonal axis, motivating the DenseRAG extension over lexical, dense SVD, neural, and cross-encoder reranking under one audit interface.

\textbf{Data valuation and influence.}
Data-valuation methods estimate how much individual records contribute to utility~\cite{ghorbani2019datashapley,jia2019data,kwon2021betashapley,wang2023databanzhaf}, and influence functions estimate the effect of training examples on predictions~\cite{koh2017influence,pruthi2020tracin,feldman2020memorization,ilyas2022datamodels}. \method{} uses this line differently: valuation prioritizes audit effort after a deletion request rather than scoring training contributions.

\textbf{Extraction, membership, canaries, and watermarking.}
Language models can expose training data through extraction~\cite{carlini2021extracting,carlini2019secretsharer,carlini2023quantifying}, membership inference~\cite{shokri2017membership,yeom2018privacy,mattern2023membership}, and canary or watermark signals~\cite{kirchenbauer2023watermark,zhao2023watermarksurvey}. These black-box channels reveal data dependence even when accuracy looks acceptable, but each is incomplete: canary-only checks fail once the marker is suppressed, membership tests miss retrieval-only leakage, and extraction prompts miss ordinary paraphrased dependence. Our benchmark therefore includes membership-inference, extraction, retriever-only, and influence-proxy baselines beside full \method{}, and incorporates triggered-leakage, graph-pivot, retrieval-dependence, and watermark evidence rather than relying on any single axis.

\textbf{Privacy, compliance, and memory governance.}
Deletion auditing is motivated by data-erasure rights, but technical audits should not overclaim legal compliance; prior work on differential privacy, certified removal, and deletion-as-control provides context~\cite{dwork2006dp,voigt2017gdpr,guo2020certified,cohen2022deletioncontrol}. \method{} contributes measurable evidence that coexists with such mechanisms: when certified removal or DP continual release is available, it monitors pipeline drift and derivative artifacts outside the formal core. Finally, adaptive-retrieval gating~\cite{wang2024ragate} motivates two diagnostics in \bench{}: an answer-only black-box audit for deployments that hide retrieval metadata, and a retrieval-off counterfactual audit that tests whether disabling retrieval changes the residual-dependence score.
